% 论文概述 —— 负责人：Sean
\subsection{论文概述}
\label{sec:cl-overview}

本节依次概述论文的研究背景、问题定义、方法框架以及实验设置与结论，内容均依据论文原文整理。
本组对方法的分析与评价集中在第~\ref{sec:cl-method}~节。

% ---------------------------------------------------------------------
\subsubsection{研究背景}

检索增强生成（Retrieval-Augmented Generation, RAG）将外部检索得到的文档拼接进提示，用以弥补大语言模型参数知识过时或缺失的不足\upcite{lewis2020retrieval,gao2023retrieval}。
随着外部知识库不断更新，检索内容与模型在预训练阶段记住的知识可能互相矛盾，这一现象称为知识冲突。
已有研究表明，发生冲突时模型往往倾向于坚持自身的记忆而忽视检索证据\upcite{longpre2021entity,xie2023adaptive}，从而降低 RAG 系统输出的准确性。

论文将已有的冲突消解方法归为两类。
词元级方法在解码阶段调整输出词元的概率分布，例如对比有无上下文时的输出分布\upcite{yuan2024discerning,jin2024tug}；
语义级方法在文段或陈述层面对齐内部与外部知识，例如 FaithfulRAG 先让模型显式写出自身的参数知识，再与检索内容比对\upcite{zhang2025faithfulrag}。
作者认为两类方法依赖的表示较为零散，难以准确刻画事实之间的关联，因而无法精确定位发生冲突的具体事实。

% ---------------------------------------------------------------------
\subsubsection{问题定义}

记 $\mathcal{M}$ 为具有参数知识 $\mathcal{K}_p$ 的大语言模型。
给定问题 $q$，RAG 系统从外部知识库 $\mathcal{K}_e$ 中检索文档，并以问题和检索内容为条件生成回答。
若存在与 $q$ 相关的两条事实陈述 $f_p\in\mathcal{K}_p$ 与 $f_e\in\mathcal{K}_e$，且二者在事实上不一致，则称发生了知识冲突。
论文采取外部知识优先的立场，目标是在冲突发生时使模型遵循检索内容作答，同时在无冲突时不损害原有性能。
围绕这一目标，论文提出了三个技术问题：如何表示事实知识以便检测冲突，如何为给定问题检索并排序相关事实，以及如何使模型在面对冲突证据时做出可靠决策。

% ---------------------------------------------------------------------
\subsubsection{方法框架}

TruthfulRAG 由图构建、图检索和冲突消解三个模块串联而成，整体流程如图~\ref{fig:cl-pipeline}~所示。

\begin{figure}[tbp]
  \centering
  \clfig{figure/cl_method/pipeline.png}{0.92\textwidth}
  \caption{TruthfulRAG 的整体流程，依次为图构建、图检索与基于熵的冲突消解（取自论文开源仓库）}
  \label{fig:cl-pipeline}
\end{figure}

图构建模块将检索内容划分为语义连贯的文段，由模型 $\mathcal{M}$ 从每个文段中抽取三元组 $(h,r,t)$，汇总得到知识图谱 $\mathcal{G}=(\mathcal{E},\mathcal{R},\mathcal{T})$，其中 $\mathcal{E}$、$\mathcal{R}$、$\mathcal{T}$ 分别为实体、关系与三元组集合。

图检索模块先从问题中抽取关键要素，按语义相似度选出最相关的实体集合 $\mathcal{E}_{imp}$ 与关系集合 $\mathcal{R}_{imp}$，再从每个关键实体出发做两跳遍历，得到初始路径集合 $\mathcal{P}_{init}$。
每条路径按其对关键实体与关系的覆盖程度打分：
\begin{equation}
  \mathrm{Ref}(p)=\alpha\,\frac{|\{e\in p\}\cap\mathcal{E}_{imp}|}{|\mathcal{E}_{imp}|}
               +\beta\,\frac{|\{r\in p\}\cap\mathcal{R}_{imp}|}{|\mathcal{R}_{imp}|},
  \label{eq:cl-ref}
\end{equation}
得分最高的 $K$ 条路径构成 $\mathcal{P}_{super}$。每条路径被展开为路径序列、实体属性与关系属性三部分文本。

冲突消解模块以输出熵度量模型的不确定性。设生成答案长度为 $L$，第 $t$ 个位置上前 $k$ 个候选词元的概率为 $p_i^{(t)}$，则
\begin{equation}
  H = -\frac{1}{L}\sum_{t=1}^{L}\sum_{i=1}^{k} p_i^{(t)}\log_2 p_i^{(t)} .
  \label{eq:cl-entropy}
\end{equation}
分别计算只给出问题时的熵 $H_{param}$ 与给出问题和路径 $p$ 时的熵 $H_{aug}(p)$，二者之差 $\Delta H_p=H_{aug}(p)-H_{param}$ 被视为冲突的指示量。
论文认为 $\Delta H_p$ 为正说明该路径与模型的参数知识相矛盾，对模型具有纠正作用，因此保留
\begin{equation}
  \mathcal{P}_{corrective}=\{\,p\in\mathcal{P}_{super}\mid \Delta H_p>\tau\,\},
  \label{eq:cl-paperrule}
\end{equation}
并以 $\mathcal{M}(q\oplus\mathcal{P}_{corrective})$ 生成最终回答，其中 $\tau$ 为阈值，$\oplus$ 表示文本拼接。

% ---------------------------------------------------------------------
\subsubsection{实验设置与主要结论}

论文在四个数据集上评测：FaithEval\upcite{ming2025faitheval} 包含反事实与不一致的上下文；MuSiQue\upcite{trivedi2022musique} 与 SQuAD\upcite{rajpurkar2016squad} 采用 KRE\upcite{ying2024intuitive} 构造的反事实版本；RealtimeQA\upcite{kasai2023realtime} 侧重时效性冲突。
所用模型为 GPT-4o-mini、Qwen2.5-7B-Instruct 与 Mistral-7B-Instruct。
对比基线包括不使用检索的直接生成、标准 RAG、提示优化方法 KRE、解码干预方法 COIECD\upcite{yuan2024discerning} 以及 FaithfulRAG，评价指标为准确率。

论文报告的主要结论有四点。
其一，TruthfulRAG 在三个模型上均取得最高的平均准确率，在 GPT-4o-mini、Qwen2.5-7B 与 Mistral-7B 上分别为 78.8\%、78.3\% 与 81.3\%，次优的 FaithfulRAG 分别为 76.5\%、78.1\% 与 80.7\%。
其二，在 MuSiQue 与 SQuAD 的无冲突版本上，该方法的准确率不低于标准 RAG。
其三，与自然语言上下文相比，结构化的推理路径使模型赋予正确答案更高的对数概率。
其四，消融实验表明图模块与冲突消解模块均对性能有贡献。
